\chapter{The Lambda Calculus and Operational Semantics}\label{ch:b3:lambda-calculus}

The interpreter of the last chapter defines Moor: whatever it does is what a Moor program means. That is a fine definition for running programs and a poor one for reasoning about them, since to know what \code{x + 1} means one must read a few hundred lines of Rust. This chapter takes the opposite road. It starts from the smallest language that can compute anything, with one way to build a function and one way to call it, and then writes the meaning of a richer language as a handful of rules, each one line, from which an interpreter can be read off.

\section{Terms and substitution}\label{sec:b3:lambda-calculus:terms}

\begin{recall}
A closure is a function value that carries the variables it uses from its surrounding scope; a higher-order function takes or returns functions (Year 2 volume). Structural induction proves a property of every tree by proving it for each kind of node, assuming it for the children (Year 1 volume). In a formula $\forall x.\,\varphi$, the quantifier binds the occurrences of $x$ in $\varphi$; the others are free (Year 2 volume, predicate logic).
\end{recall}

\begin{definition}[$\lambda$-term]\label{def:b3:lambda-calculus:terms}
The \emph{$\lambda$-terms}\index{lambda-term@$\lambda$-term} over a set of variables are given by the grammar $M ::= x \mid \lambda x.\,M \mid M\,M$. A term $\lambda x.\,M$ is an \emph{abstraction}\index{abstraction (lambda-calculus)@abstraction ($\lambda$-calculus)}: the function that maps $x$ to $M$. A term $M\,N$ is an \emph{application}\index{application (lambda-calculus)@application ($\lambda$-calculus)}: the function $M$ called with the argument $N$.
\end{definition}

\omterm{def:b3:lambda-calculus:terms}{Application} associates to the left, so $f\,a\,b$ is $(f\,a)\,b$, a function of two arguments taken one at a time; and the body of a $\lambda$ extends as far right as possible, so $\lambda x.\,\lambda y.\,x\,y$ is $\lambda x.\,(\lambda y.\,(x\,y))$. In code we write \texttt{\textbackslash{}x.\ M} for $\lambda x.\,M$, and \texttt{\textbackslash{}x y.\ M} for $\lambda x.\,\lambda y.\,M$. A $\lambda$ binds its variable in its body exactly as a quantifier does in logic: in $\lambda x.\,x\,y$ the occurrence of $x$ is bound and $y$ is free, and the free variables of a term are defined by the same structural recursion as for formulas.

\begin{definition}[$\alpha$-equivalence, capture-avoiding substitution]\label{def:b3:lambda-calculus:alpha}
Two terms are \emph{$\alpha$-equivalent}\index{alpha-equivalence@$\alpha$-equivalence} if they differ only in the names of bound variables, renamed consistently: $\lambda x.\,x\,y$ and $\lambda z.\,z\,y$ are, $\lambda x.\,x\,y$ and $\lambda y.\,y\,y$ are not. The \emph{capture-avoiding substitution}\index{capture-avoiding substitution} $M[x := N]$ replaces the free occurrences of $x$ in $M$ by $N$, after renaming the bound variables of $M$ that would otherwise capture a free variable of $N$.
\end{definition}

Capture is the whole difficulty. Substituting $y$ for $x$ in $\lambda y.\,x$ naively gives $\lambda y.\,y$, the identity, where we wanted the function that ignores its argument and returns $y$. The correct result renames first: $\lambda y'.\,y$. Terms are considered up to $\alpha$-equivalence from now on, so this renaming is harmless.

Names make implementations fragile: every substitution must check for capture and invent fresh names. Nicolaas \omterm{def:b3:lambda-calculus:debruijn}{de Bruijn} removed the names altogether.

\begin{definition}[de Bruijn index]\label{def:b3:lambda-calculus:debruijn}
In the \emph{de Bruijn}\index{de Bruijn index} notation, a bound variable is replaced by the number of $\lambda$'s between its occurrence and its binder: 0 for the nearest enclosing $\lambda$, 1 for the next one out, and so on. The term $\lambda x.\,\lambda y.\,x\,y$ becomes $\lambda.\,\lambda.\,1\;0$.
\end{definition}

Two $\alpha$-equivalent terms have the same \omterm{def:b3:lambda-calculus:debruijn}{de Bruijn} form, so $\alpha$-equivalence becomes plain equality of trees. \Cref{lst:b3:lambda-calculus:term} is the representation used in this chapter: an index for a bound variable, a name for a free one, and the binder's name kept only as a hint for printing.

\omcode{code/bachelor-3/05-lambda-calculus/listings/term.rs}{8}{33}{$\lambda$-terms with \omterm{def:b3:lambda-calculus:debruijn}{de Bruijn} indices; $\alpha$-equivalence is structural equality, with the names ignored.}\label{lst:b3:lambda-calculus:term}

\begin{omfigure}
\begin{tikzpicture}[font=\small, >=Stealth,
  nd/.style={inner sep=2pt, minimum size=5mm},
  bnd/.style={->, dashed, omThm, thick}]
  \node[nd] (lx) at (0,0) {$\lambda x$};
  \node[nd] (ly) at (0,-0.9) {$\lambda y$};
  \node[nd] (ap) at (0,-1.8) {@};
  \node[nd] (x) at (-1.3,-2.7) {$x$ : 1};
  \node[nd] (lz) at (1.3,-2.7) {$\lambda z$};
  \node[nd] (ap2) at (1.3,-3.6) {@};
  \node[nd] (z) at (0.5,-4.5) {$z$ : 0};
  \node[nd] (y) at (2.1,-4.5) {$y$ : 1};
  \draw (lx) -- (ly) -- (ap) -- (x); \draw (ap) -- (lz) -- (ap2) -- (z); \draw (ap2) -- (y);
  \draw[bnd] (x.west) to[bend left=40] (lx.west);
  \draw[bnd] (z.west) to[bend left=50] (lz.west);
  \draw[bnd] (y.east) to[bend right=40] (ly.east);
  \node[anchor=west, align=left] at (3.6,-1.6) {named: $\lambda x.\,\lambda y.\,x\,(\lambda z.\,z\,y)$\\[2pt] de Bruijn: $\lambda.\,\lambda.\,1\,(\lambda.\,0\;1)$};
\end{tikzpicture}
\par\smallskip{\small A term as a tree, each variable joined to its binder (dashed). The index counts the $\lambda$ nodes crossed on the way up: from $y$ the path meets $\lambda z$ first, so $y$ is 1, and so is $x$, which meets $\lambda y$.}
\end{omfigure}

Substitution must now keep indices right. When a term $N$ is substituted under a $\lambda$, its free indices must grow by one, since they now cross one more binder: this is the \emph{shift}. And when a $\lambda$ disappears in a reduction, the indices that pointed past it must shrink by one. \Cref{lst:b3:lambda-calculus:subst} does both.

\omcode{code/bachelor-3/05-lambda-calculus/listings/reduce.rs}{6}{31}{Shifting and substitution on \omterm{def:b3:lambda-calculus:debruijn}{de Bruijn} terms. No name is ever compared, so no capture can happen.}\label{lst:b3:lambda-calculus:subst}

\section{Reduction}\label{sec:b3:lambda-calculus:reduction}

\begin{definition}[$\beta$-reduction, redex, normal form]\label{def:b3:lambda-calculus:beta}
A \emph{redex}\index{redex} is a subterm of the form $(\lambda x.\,M)\,N$. \emph{$\beta$-reduction}\index{beta-reduction@$\beta$-reduction} replaces one redex by $M[x := N]$, anywhere in a term; we write $M \to_\beta M'$ for one step and $\to_\beta^*$ for zero or more. A term with no redex is in \emph{normal form}\index{normal form (of a lambda-term)@normal form (of a $\lambda$-term)}.
\end{definition}

The term $(\lambda x.\,\lambda y.\,x\,y)\,y$ reduces in one step to $\lambda y'.\,y\,y'$, renaming to avoid capture. Some terms have no \omterm{def:b3:lambda-calculus:beta}{normal form}: $\Omega = (\lambda x.\,x\,x)\,(\lambda x.\,x\,x)$ reduces to itself, forever. And a term with several \omterm{def:b3:lambda-calculus:beta}{redexes} may be reduced in several orders, which is where the choices of a programming language begin.

\begin{definition}[Evaluation strategy]\label{def:b3:lambda-calculus:strategy}
An \emph{evaluation strategy}\index{evaluation strategy} chooses which \omterm{def:b3:lambda-calculus:beta}{redex} to reduce next. \emph{Normal-order reduction}\index{normal-order reduction} reduces the leftmost outermost \omterm{def:b3:lambda-calculus:beta}{redex}, also under a $\lambda$. \emph{Call by name}\index{call by name} does the same but never reduces under a $\lambda$ or inside an argument. \emph{Call by value}\index{call by value} reduces a \omterm{def:b3:lambda-calculus:beta}{redex} $(\lambda x.\,M)\,N$ only once $N$ is a value (here, an \omterm{def:b3:lambda-calculus:terms}{abstraction}), and never under a $\lambda$. \emph{Lazy evaluation}\index{lazy evaluation} (call by need) is call by name in which an argument, once evaluated, is shared by all the places that use it.
\end{definition}

\Cref{lst:b3:lambda-calculus:step} implements four strategies: the three above and \emph{applicative order}, leftmost innermost, which reduces arguments first and also goes under $\lambda$'s; it is \omterm{def:b3:lambda-calculus:strategy}{call by value} continued to a \omterm{def:b3:lambda-calculus:beta}{normal form}. Each step returns the new term, or nothing when the strategy finds no \omterm{def:b3:lambda-calculus:beta}{redex}.

\omcode{code/bachelor-3/05-lambda-calculus/listings/reduce.rs}{46}{78}{One step of each strategy. Normal order and \omterm{def:b3:lambda-calculus:strategy}{call by name} contract a \omterm{def:b3:lambda-calculus:beta}{redex} at the top before looking inside; \omterm{def:b3:lambda-calculus:strategy}{call by value} and applicative order look inside first.}\label{lst:b3:lambda-calculus:step}

The reduction graph below shows every way to reduce one small term. Each strategy is one path through it; they differ in length, and \omterm{def:b3:lambda-calculus:strategy}{call by name} even stops early, at a term whose head is the free variable $a$, since it never reduces inside an argument.

\begin{omfigure}
\begin{tikzpicture}[font=\small, >=Stealth, node distance=8mm,
  t/.style={draw, rounded corners=2pt, fill=omDef!6, inner sep=3pt}]
  \node[t] (T) at (0,0) {$(\lambda x.\,x\,x)\,((\lambda y.\,y)\,a)$};
  \node[t] (A) at (-3.3,-1.5) {$(\lambda y.\,y)\,a\,((\lambda y.\,y)\,a)$};
  \node[t] (B) at (3.3,-1.5) {$(\lambda x.\,x\,x)\,a$};
  \node[t] (C) at (-5.0,-3.0) {$a\,((\lambda y.\,y)\,a)$};
  \node[t] (D) at (-1.5,-3.0) {$(\lambda y.\,y)\,a\,a$};
  \node[t] (N) at (0,-4.5) {$a\,a$};
  \draw[->, very thick, omThm] (T) -- (A);
  \draw[->, very thick, omMeth] (T) -- (B);
  \draw[->, very thick, omThm] (A) -- (C);
  \draw[->] (A) -- (D);
  \draw[->, very thick, omMeth] (B) -- (N);
  \draw[->, very thick, omThm] (C) -- (N);
  \draw[->] (D) -- (N);
  \node[omThm, anchor=west, align=left, font=\footnotesize] at (-7.2,-4.4) {normal order: 3 steps\\call by name: the first 2, then stops};
  \node[omMeth, anchor=west, align=left, font=\footnotesize] at (2.0,-3.4) {call by value and\\applicative order: 2 steps};
\end{tikzpicture}
\par\smallskip{\small All the reducts of one term, an arrow per $\beta$-step (\code{reduction_graph_of_the_figure} enumerates them: six terms, seven arrows). Normal order copies the unreduced argument and reduces it twice; \omterm{def:b3:lambda-calculus:strategy}{call by value} reduces it once, before the copy.}
\end{omfigure}

Every path that ends, ends at $a\,a$. This is no accident of the example. To say why, we step back to rewriting in general.

\begin{definition}[Term rewriting system, confluence]\label{def:b3:lambda-calculus:rewriting}
A \emph{rewriting system}\index{term rewriting system} is a set $A$ with a relation $\to$ on it, often given by rules that rewrite terms (a \emph{term rewriting system}); $\to^*$ is its reflexive transitive closure. The system is \emph{confluent}\index{confluence} if whenever $a \to^* b$ and $a \to^* c$, there is a $d$ with $b \to^* d$ and $c \to^* d$; \emph{locally confluent} if this holds whenever $a \to b$ and $a \to c$ in one step each; \emph{terminating} if there is no infinite chain $a_0 \to a_1 \to \cdots$.
\end{definition}

\begin{theorem}[Church--Rosser]\label{thm:b3:lambda-calculus:church-rosser}
$\beta$-reduction is \omterm{def:b3:lambda-calculus:rewriting}{confluent}.\index{Church--Rosser theorem}
\end{theorem}

\admitted

Local \omterm{def:b3:lambda-calculus:rewriting}{confluence} is not enough here, since $\beta$-reduction does not terminate, and a one-step diamond fails: reducing a \omterm{def:b3:lambda-calculus:beta}{redex} inside an argument that is copied twice must be repeated in both copies. The usual proof defines a \emph{parallel reduction} that contracts any set of \omterm{def:b3:lambda-calculus:beta}{redexes} at once, shows by induction on terms that it satisfies the one-step diamond, and notes that its reflexive transitive closure is $\to_\beta^*$.

\begin{corollary}\label{cor:b3:lambda-calculus:unique}
A $\lambda$-term has at most one \omterm{def:b3:lambda-calculus:beta}{normal form}, up to $\alpha$-equivalence.
\end{corollary}

\begin{proof}
If $M \to_\beta^* N_1$ and $M \to_\beta^* N_2$ with both in \omterm{def:b3:lambda-calculus:beta}{normal form}, \omterm{def:b3:lambda-calculus:rewriting}{confluence} gives $P$ with $N_1 \to_\beta^* P$ and $N_2 \to_\beta^* P$. A \omterm{def:b3:lambda-calculus:beta}{normal form} reduces only to itself, so $N_1 = P = N_2$.
\end{proof}

So the strategies cannot disagree on the result; they can only disagree on whether they reach it. Normal order is the safe one.

\begin{theorem}[Normal order normalises]\label{thm:b3:lambda-calculus:normal-order}
If a $\lambda$-term has a \omterm{def:b3:lambda-calculus:beta}{normal form}, \omterm{def:b3:lambda-calculus:strategy}{normal-order reduction} reaches it.
\end{theorem}

\admitted

The proof goes through the \emph{standardisation theorem}: any reduction to a \omterm{def:b3:lambda-calculus:beta}{normal form} can be reordered so that it contracts \omterm{def:b3:lambda-calculus:beta}{redexes} from left to right, outermost first. It is longer than this chapter allows; we use the result and check it on examples. Applicative order is the fragile one: in $(\lambda x.\,\lambda y.\,y)\,\Omega$ it reduces $\Omega$ forever, while normal order discards it in one step (exercise~\ref{exo:b3:lambda-calculus:11}).

For a terminating system, local \omterm{def:b3:lambda-calculus:rewriting}{confluence} is all one needs to check, because of a lemma whose proof is a model of well-founded induction.

\begin{lemma}[Newman]\label{lem:b3:lambda-calculus:newman}
A terminating, locally \omterm{def:b3:lambda-calculus:rewriting}{confluent} \omterm{def:b3:lambda-calculus:rewriting}{rewriting system} is \omterm{def:b3:lambda-calculus:rewriting}{confluent}.\index{Newman's lemma}
\end{lemma}

\begin{proof}
Since the system terminates, the relation ``$b$ is reachable from $a$ in at least one step'' has no infinite chain, so we may prove a property of all elements by induction along it: assume the property for every element reachable from $a$ in one or more steps, and prove it for $a$. The property is: every two reducts of $a$ have a common reduct.

Let $a \to^* b$ and $a \to^* c$. If either reduction is empty, say $a = b$, take $d = c$. Otherwise write $a \to b_1 \to^* b$ and $a \to c_1 \to^* c$. By local \omterm{def:b3:lambda-calculus:rewriting}{confluence}, there is $e$ with $b_1 \to^* e$ and $c_1 \to^* e$. Now $b_1$ is a reduct of $a$, so the induction hypothesis applies to it: $b$ and $e$, both reducts of $b_1$, have a common reduct $f$. Likewise $c_1$ is a reduct of $a$, and $f$ and $c$ are both reducts of $c_1$ (since $c_1 \to^* e \to^* f$), so they have a common reduct $d$. Then $b \to^* f \to^* d$ and $c \to^* d$.
\end{proof}

\begin{omfigure}
\begin{tikzpicture}[font=\small, >=Stealth, x=1.15cm, y=1.0cm]
  \begin{scope}
  \node (a) at (0,0) {$a$}; \node (b) at (-1.2,-1.2) {$b$}; \node (c) at (1.2,-1.2) {$c$}; \node (d) at (0,-2.4) {$d$};
  \draw[->] (a) -- node[above left] {$*$} (b); \draw[->] (a) -- node[above right] {$*$} (c);
  \draw[->, dashed] (b) -- node[below left] {$*$} (d); \draw[->, dashed] (c) -- node[below right] {$*$} (d);
  \node at (0,-3.2) {confluence};
  \end{scope}
  \begin{scope}[xshift=5.2cm]
  \node (a) at (0,0) {$a$};
  \node (b1) at (-1.1,-1) {$b_1$}; \node (c1) at (1.1,-1) {$c_1$};
  \node (b) at (-2.2,-2) {$b$}; \node (e) at (0,-2) {$e$}; \node (c) at (2.2,-3) {$c$};
  \node (f) at (-1.1,-3) {$f$}; \node (d) at (0.55,-4) {$d$};
  \draw[->] (a) -- (b1); \draw[->] (a) -- (c1);
  \draw[->] (b1) -- node[above left] {$*$} (b); \draw[->] (c1) -- node[right, pos=0.3] {$*$} (c);
  \draw[->, dashed, omMeth] (b1) -- node[below left, pos=0.7] {$*$} (e); \draw[->, dashed, omMeth] (c1) -- node[below right, pos=0.7] {$*$} (e);
  \draw[->, dashed, omThm] (b) -- (f); \draw[->, dashed, omThm] (e) -- (f);
  \draw[->, dashed, omDef] (f) -- (d); \draw[->, dashed, omDef] (c) -- (d);
  \node[omMeth, font=\footnotesize] at (0,-1.15) {local};
  \node[omThm, font=\footnotesize, align=center] at (-2.4,-3.0) {induction\\at $b_1$};
  \node[omDef, font=\footnotesize, align=center] at (2.3,-4.0) {induction\\at $c_1$};
  \end{scope}
\end{tikzpicture}
\par\smallskip{\small Left: what \omterm{def:b3:lambda-calculus:rewriting}{confluence} asks. Right: the proof of Newman's lemma, three tiles. Local \omterm{def:b3:lambda-calculus:rewriting}{confluence} closes the top diamond; the induction hypothesis, used at the two one-step reducts $b_1$ and $c_1$, closes the two others. Termination is what makes the induction legitimate.}
\end{omfigure}

Termination is usually proved by a \emph{measure}: a function to the natural numbers that every step decreases. The two rules $\mathit{add}(0, y) \to y$ and $\mathit{add}(s(x), y) \to s(\mathit{add}(x, y))$ define addition on numbers written with $0$ and a successor $s$. Give each term the weight $w$ with $w(0) = 1$, $w(s(t)) = w(t) + 1$, $w(\mathit{add}(t, u)) = 2w(t) + w(u)$: the first rule removes at least $w(y) + 2 - w(y) = 2$, the second $2(w(x)+1) + w(y) - 2w(x) - w(y) - 1 = 1$, and weights are positive, so every chain is finite. Two rules can only conflict where one's left side overlaps another's; here none does, so the system is locally \omterm{def:b3:lambda-calculus:rewriting}{confluent}, and \omterm{def:b3:lambda-calculus:rewriting}{confluent} by Newman's lemma: every closed sum has exactly one value. Without termination the lemma fails: exercise~\ref{exo:b3:lambda-calculus:10} gives four elements that show it.

\section{Data as functions}\label{sec:b3:lambda-calculus:church}

The $\lambda$-calculus has no numbers, no booleans and no pairs, and needs none: each can be represented by what it does.

\begin{definition}[Church boolean, Church numeral]\label{def:b3:lambda-calculus:church}
The \emph{Church booleans}\index{Church boolean} are $\mathit{true} = \lambda a.\,\lambda b.\,a$ and $\mathit{false} = \lambda a.\,\lambda b.\,b$: a boolean chooses between two arguments. The \emph{Church numeral}\index{Church numeral} of $n$ is $c_n = \lambda f.\,\lambda x.\,f\,(f\,(\cdots(f\,x)))$, with $n$ \omterm{def:b3:lambda-calculus:terms}{applications} of $f$: the number $n$ is the operation ``do something $n$ times''.
\end{definition}

A conditional needs no keyword: $b\,M\,N$ is $M$ when $b$ is $\mathit{true}$ and $N$ when it is $\mathit{false}$. A pair holds its two components until a selector is given to it: $\mathit{pair} = \lambda a.\,\lambda b.\,\lambda f.\,f\,a\,b$, with $\mathit{fst}\,p = p\,\mathit{true}$. Arithmetic follows from ``$n$ times'': $\mathit{succ}\,n$ does $f$ once more than $n$ does; $\mathit{plus}\,m\,n$ does $f$ $m$ times after $n$ times; $\mathit{mult}\,m\,n$ does ``$f$ $n$ times'', $m$ times; and $\mathit{exp}\,m\,n = n\,m$ applies ``multiply by $m$'' $n$ times, starting from the function itself.

\omcode{code/bachelor-3/05-lambda-calculus/listings/encode.rs}{21}{36}{Booleans, pairs and arithmetic as $\lambda$-terms. Each name may use the earlier ones; \code{parse_with} replaces a name by its term.}\label{lst:b3:lambda-calculus:encode}

The predecessor is the hard one, because a numeral can only do something $n$ times, never undo a step. The trick is to count pairs: start from $(0, 0)$ and apply $n$ times the step $(a, b) \mapsto (b, b + 1)$. After $k$ steps the pair is $(k - 1, k)$ for $k \ge 1$, so after $n$ steps its first component is $n - 1$ (and $0$ when $n = 0$). The definition \code{pred} in \cref{lst:b3:lambda-calculus:encode} is exactly that.

\omoutput{code/bachelor-3/05-lambda-calculus/outputs/church.txt}{Six computations, each reduced to \omterm{def:b3:lambda-calculus:beta}{normal form} under normal order and under applicative order, with the number of $\beta$-steps. The \omterm{def:b3:lambda-calculus:beta}{normal forms} are recognised as numerals; \code{exp c2 c3} prints its numeral with the names \code{x} and \code{x'}, an $\alpha$-variant of $c_8$.}\label{out:b3:lambda-calculus:church}

Addition, multiplication and exponentiation each take a handful of steps, the same under both strategies; the predecessor takes dozens. The weekend problem measures how the counts grow with the numbers.

\section{Recursion without names}\label{sec:b3:lambda-calculus:fixpoint}

The $\lambda$-calculus has no way to name a function, so a function cannot call itself by name. It can, however, receive itself as an argument. Write the factorial as a functional that takes ``the function to call recursively'' as a parameter:
\[ F = \lambda \mathit{fact}.\,\lambda n.\,\mathtt{if}\ n = 0\ \mathtt{then}\ 1\ \mathtt{else}\ n \times \mathit{fact}\,(n - 1). \]
The factorial is the function $g$ with $F\,g = g$: a fixed point of $F$.

\begin{definition}[Fixed-point combinator]\label{def:b3:lambda-calculus:fixpoint}
A \emph{fixed-point combinator}\index{fixed-point combinator} is a closed term $\Theta$ such that $\Theta\,F =_\beta F\,(\Theta\,F)$ for every term $F$, where $=_\beta$ is the equivalence generated by $\beta$-reduction.
\end{definition}

The best known is \[ Y = \lambda f.\,(\lambda x.\,f\,(x\,x))\,(\lambda x.\,f\,(x\,x)). \] Indeed $Y\,F \to_\beta (\lambda x.\,F\,(x\,x))\,(\lambda x.\,F\,(x\,x)) \to_\beta F\,((\lambda x.\,F\,(x\,x))\,(\lambda x.\,F\,(x\,x)))$, and the argument of $F$ is the term $Y\,F$ reduced to in one step. Under \omterm{def:b3:lambda-calculus:strategy}{call by value}, $Y$ is useless: the argument $x\,x$ is evaluated before $F$ is ever entered, and it unfolds again, without end. Wrapping it in a $\lambda$ delays it: \[ Z = \lambda f.\,(\lambda x.\,f\,(\lambda v.\,x\,x\,v))\,(\lambda x.\,f\,(\lambda v.\,x\,x\,v)) \] is a \omterm{def:b3:lambda-calculus:fixpoint}{fixed-point combinator} for \omterm{def:b3:lambda-calculus:strategy}{call by value}, since $\lambda v.\,x\,x\,v$ is already a value. \Cref{lst:b3:lambda-calculus:ycomb} runs both, in the small language of the next section, which has integers and \code{if}.

\omcode{code/bachelor-3/05-lambda-calculus/listings/ycomb.rs}{9}{32}{The factorial of 5 three ways: through $Z$, through \code{letrec}, and through $Y$, all under \omterm{def:b3:lambda-calculus:strategy}{call by value}.}\label{lst:b3:lambda-calculus:ycomb}

\omoutput{code/bachelor-3/05-lambda-calculus/outputs/ycomb.txt}{$Z$ computes 120 in 46 steps, \code{letrec} in 35; $Y$ never reaches a value, and the term it rewrites keeps growing.}\label{out:b3:lambda-calculus:ycomb}

A language does not need $Z$ in its programs: \texttt{letrec f = \textbackslash{}x.\ e1 in e2} is sugar for binding $f$ to the fixed point of $\lambda f.\,\lambda x.\,e_1$, and the semantics of the next section gives it directly, by one rule that unfolds the definition once. It saves the eleven steps $Z$ spends rebuilding itself.

\section{Small steps and big steps}\label{sec:b3:lambda-calculus:semantics}

The $\lambda$-language of this chapter and the next, \code{.lam} files, adds to the $\lambda$-calculus the integers, the booleans, \code{let}, \code{letrec}, \code{if}, pairs and five operators. Its syntax tree is in \cref{lst:b3:lambda-calculus:lam}; the optional type on a $\lambda$ is for chapter~\ref{ch:b3:type-systems}.

\omcode{code/bachelor-3/05-lambda-calculus/listings/lam.rs}{13}{26}{The expressions of the $\lambda$-language.}\label{lst:b3:lambda-calculus:lam}

\begin{definition}[Operational semantics, small-step, big-step]\label{def:b3:lambda-calculus:semantics}
An \emph{operational semantics}\index{operational semantics} defines the meaning of a program by rules describing how it runs. A \emph{small-step semantics}\index{small-step semantics} defines a relation $e \to e'$, one step of computation, and a program's result is the value $v$ with $e \to^* v$. A \emph{big-step semantics}\index{big-step semantics} defines a relation $e \Downarrow v$, ``$e$ evaluates to $v$'', directly, by rules whose premises evaluate the parts of $e$.
\end{definition}

The values are the integers, the booleans, the \omterm{def:b3:lambda-calculus:terms}{abstractions}, and the pairs of values. The program is closed, and \omterm{def:b3:lambda-calculus:strategy}{call by value} only ever substitutes values, which are then closed too: nothing can be captured, and substitution by name is safe. These are the small-step rules for \omterm{def:b3:lambda-calculus:terms}{application}, \code{let}, \code{letrec} and \code{if}; an operator applied to two values computes ($2 + 3 \to 5$), and its operands are reduced left to right like an \omterm{def:b3:lambda-calculus:terms}{application}'s.
\begin{mathpar}
\inferrule{ }{(\lambda x.\,e)\,v \to e[x := v]}

\inferrule{e_1 \to e_1'}{e_1\,e_2 \to e_1'\,e_2}

\inferrule{e_2 \to e_2'}{v\,e_2 \to v\,e_2'}

\inferrule{ }{\mathtt{let}\ x = v\ \mathtt{in}\ e \to e[x := v]}

\inferrule{ }{\mathtt{letrec}\ f = \lambda x.\,e_1\ \mathtt{in}\ e_2 \to e_2[f := \lambda x.\,\mathtt{letrec}\ f = \lambda x.\,e_1\ \mathtt{in}\ e_1]}

\inferrule{ }{\mathtt{if}\ \mathtt{true}\ \mathtt{then}\ e_2\ \mathtt{else}\ e_3 \to e_2}

\inferrule{e_1 \to e_1'}{\mathtt{if}\ e_1\ \mathtt{then}\ e_2\ \mathtt{else}\ e_3 \to \mathtt{if}\ e_1'\ \mathtt{then}\ e_2\ \mathtt{else}\ e_3}
\end{mathpar}
The rules without premises \emph{compute}; those with a premise only say where to look. The interpreter of \cref{lst:b3:lambda-calculus:small} is read off them, one match arm per rule, in the same order; a term no rule applies to is stuck, and the interpreter says why.

\omcode{code/bachelor-3/05-lambda-calculus/listings/smallstep.rs}{55}{89}{The small-step rules of the $\lambda$-language as a function: one arm per rule.}\label{lst:b3:lambda-calculus:small}

The big-step rules say the same thing at once, from the leaves of a derivation up:
\begin{mathpar}
\inferrule{ }{v \Downarrow v}

\inferrule{e_1 \Downarrow \lambda x.\,e \\ e_2 \Downarrow v_2 \\ e[x := v_2] \Downarrow v}{e_1\,e_2 \Downarrow v}

\inferrule{e_1 \Downarrow v_1 \\ e_2[x := v_1] \Downarrow v}{\mathtt{let}\ x = e_1\ \mathtt{in}\ e_2 \Downarrow v}

\inferrule{e_1 \Downarrow \mathtt{true} \\ e_2 \Downarrow v}{\mathtt{if}\ e_1\ \mathtt{then}\ e_2\ \mathtt{else}\ e_3 \Downarrow v}

\inferrule{e_1 \Downarrow n_1 \\ e_2 \Downarrow n_2}{e_1 + e_2 \Downarrow n_1 + n_2}
\end{mathpar}
Read as an algorithm, the big-step rules are a recursive evaluator, the shape of chapter~\ref{ch:b3:semantic-analysis}'s interpreter: each recursive call builds one premise (\cref{lst:b3:lambda-calculus:big}).

\omcode{code/bachelor-3/05-lambda-calculus/listings/bigstep.rs}{21}{53}{The big-step rules as an evaluator. Each call is one node of the derivation; \code{z} counts the nodes, and those whose rule computes.}\label{lst:b3:lambda-calculus:big}

\begin{omfigure}
\begin{mathpar}
\inferrule*[right=app]{\inferrule*{ }{\lambda x.\,x + 1 \Downarrow \lambda x.\,x + 1} \\ \inferrule*{ }{2 \Downarrow 2} \\ \inferrule*[right=add]{\inferrule*{ }{2 \Downarrow 2} \\ \inferrule*{ }{1 \Downarrow 1}}{2 + 1 \Downarrow 3}}{(\lambda x.\,x + 1)\,2 \Downarrow 3}
\end{mathpar}
\par\smallskip{\small The derivation of $(\lambda x.\,x + 1)\,2 \Downarrow 3$: six nodes, two of which compute (the \omterm{def:b3:lambda-calculus:terms}{application} and the addition). The small-step run takes exactly two steps: $(\lambda x.\,x + 1)\,2 \to 2 + 1 \to 3$.}
\end{omfigure}

\Cref{out:b3:lambda-calculus:sum} runs \code{sum.lam}, $3 + 2 + 1 + 0$ by recursion, both ways. The small-step trace shows what the big-step derivation hides: the pending additions pile up on the left, \code{3 + (2 + (1 + ...))}, and are only performed once the recursion has bottomed out. The demonstration prints the unfolded function as \code{<sum>} and an unfolded \code{letrec} with its definition elided.

\omoutput{code/bachelor-3/05-lambda-calculus/outputs/semantics_demo.txt}{\code{sum.lam} under the small-step rules, one line per step, then under the big-step rules: 23 steps, and a derivation with 23 computing nodes.}\label{out:b3:lambda-calculus:sum}

The two counts are equal, and this is a theorem.

\begin{theorem}[Small steps and big steps agree]\label{thm:b3:lambda-calculus:agree}
For a closed expression $e$ and a value $v$: $e \to^* v$ if and only if $e \Downarrow v$. Moreover, the number of steps equals the number of computing nodes of the derivation.
\end{theorem}

\begin{proof}
($\Leftarrow$) By induction on the derivation. A value takes zero steps. For an \omterm{def:b3:lambda-calculus:terms}{application}, the premises give, by induction, $e_1 \to^* \lambda x.\,e$ in $k_1$ steps, $e_2 \to^* v_2$ in $k_2$ and $e[x := v_2] \to^* v$ in $k_3$. The congruence rules lift the first two reductions to the whole term (a short induction on their length): $e_1\,e_2 \to^* (\lambda x.\,e)\,e_2 \to^* (\lambda x.\,e)\,v_2$, then one computing step gives $e[x := v_2]$, and $k_3$ more give $v$. In all $k_1 + k_2 + 1 + k_3$ steps, the node itself being the one computing step. The other rules are alike.

($\Rightarrow$) First a lemma: if $e \to e'$ and $e' \Downarrow v$ by a derivation with $k$ computing nodes, then $e \Downarrow v$ with $k + 1$. By induction on the step. If it is $(\lambda x.\,e_0)\,v_2 \to e_0[x := v_2]$, build the \omterm{def:b3:lambda-calculus:terms}{application} node from $\lambda x.\,e_0 \Downarrow \lambda x.\,e_0$, $v_2 \Downarrow v_2$ and the given derivation. If it is $e_1\,e_2 \to e_1'\,e_2$ with $e_1 \to e_1'$, the derivation of $e_1'\,e_2 \Downarrow v$ contains $e_1' \Downarrow \lambda x.\,e_0$; by induction $e_1 \Downarrow \lambda x.\,e_0$ with one more computing node; put it back in place. The other rules are alike. Then the theorem follows by induction on the number of steps of $e \to^* v$, starting from $v \Downarrow v$.
\end{proof}

The test \code{small_and_big_step_agree} checks the theorem on six programs. It says nothing of programs that do not terminate: a small-step run of $\Omega$ is an infinite sequence, while $\Omega$ has no derivation at all. \omterm{def:b3:lambda-calculus:semantics}{Big-step semantics} cannot tell a loop from a stuck program, both having no derivation; \omterm{def:b3:lambda-calculus:semantics}{small-step semantics} can, which is why proofs of type safety in the next chapter use it.

\begin{remark}
A third style, \emph{denotational semantics}, gives each program a mathematical object as its meaning: a function from inputs to outputs, built compositionally from the meanings of the parts. A recursive definition then denotes a fixed point of a function on a suitable ordered set, and a loop denotes an ``undefined'' element. It answers questions \omterm{def:b3:lambda-calculus:semantics}{operational semantics} answers badly, such as when two programs are interchangeable in every context; it needs more mathematics than this chapter has room for.
\end{remark}

\begin{inpractice}[Keep the rules runnable]
Rules written on paper drift away from the compiler they describe. Rules read off as an interpreter, as here, can be run on the compiler's test programs: a disagreement between the two is then a bug in one of them, and the rules, being short, are the easier to check. The compiler of chapter~\ref{ch:b3:code-generation} is held to the interpreter of chapter~\ref{ch:b3:semantic-analysis} in this way.
\end{inpractice}

\begin{history}[Functions as the foundation]
Alonzo Church introduced the $\lambda$-notation in 1932, in a system meant as a foundation of logic. % ledger: b3:church-1932
In 1936 he and J. Barkley Rosser proved the \omterm{def:b3:lambda-calculus:rewriting}{confluence} theorem that bears their names. % ledger: b3:church-rosser-1936
Max Newman's lemma appeared in 1942, in a paper on equivalence defined by combinatorial rules. % ledger: b3:newman-1942
Nicolaas \omterm{def:b3:lambda-calculus:debruijn}{de Bruijn}'s nameless notation was published in 1972, as a tool for manipulating formulas by machine. % ledger: b3:debruijn-1972
The big-step style of semantics was named \emph{natural semantics} by Gilles Kahn in 1987. % ledger: b3:kahn-1987
\end{history}

\section{Exercises}\label{sec:b3:lambda-calculus:exercises}

\begin{exercise}[$\star$]\label{exo:b3:lambda-calculus:1}
Give the free variables of \texttt{\textbackslash{}x.\ x y}, of \texttt{(\textbackslash{}x.\ x) x} and of \texttt{\textbackslash{}x.\ \textbackslash{}y.\ z (x y) w}.
\end{exercise}

\begin{exercise}[$\star$]\label{exo:b3:lambda-calculus:2}
Which of these pairs are $\alpha$-equivalent? \texttt{\textbackslash{}x.\ \textbackslash{}y.\ x y} and \texttt{\textbackslash{}y.\ \textbackslash{}x.\ y x}; \texttt{\textbackslash{}x.\ \textbackslash{}y.\ x y} and \texttt{\textbackslash{}x.\ \textbackslash{}y.\ y x}; \texttt{\textbackslash{}x.\ x z} and \texttt{\textbackslash{}z.\ z z}. Give the \omterm{def:b3:lambda-calculus:debruijn}{de Bruijn} form of each term.
\end{exercise}

\begin{exercise}[$\star$]\label{exo:b3:lambda-calculus:3}
Reduce \texttt{(\textbackslash{}x.\ \textbackslash{}y.\ x y) y} by one $\beta$-step. What does a substitution that ignores capture give, and why is it wrong?
\end{exercise}

\begin{exercise}[$\star$]\label{exo:b3:lambda-calculus:4}
Predict the lines that \code{church.rs} would print for \code{succ c4} and for \code{plus c4 c0}: both step counts and the numeral.
\end{exercise}

\begin{exercise}[$\star\star$]\label{exo:b3:lambda-calculus:5}
A first version of \code{subst} handled \omterm{def:b3:lambda-calculus:terms}{abstractions} with \code{Lam(x, b) => Lam(x.clone(), Box::new(subst(j + 1, s, b)))}. Find a term that normal order reduces wrongly with it, say what it gives, and fix the line.
\end{exercise}

\begin{exercise}[$\star\star$]\label{exo:b3:lambda-calculus:6}
With the definitions of \cref{lst:b3:lambda-calculus:encode} and $\mathit{sub} = \lambda m.\,\lambda n.\,n\ \mathit{pred}\ m$, which computes $\max(m - n, 0)$, write $\mathit{is\_zero}$, which maps $c_0$ to $\mathit{true}$ and every other numeral to $\mathit{false}$, and $\mathit{eq}$, which tests two numerals for equality. Test them on all numerals up to 5.
\end{exercise}

\begin{exercise}[$\star\star$]\label{exo:b3:lambda-calculus:7}
Prove that if two strategies both lead a term to a \omterm{def:b3:lambda-calculus:beta}{normal form}, they lead it to the same one. Then reduce \code{plus c2 c3} by \omterm{def:b3:lambda-calculus:strategy}{call by name}: where does it stop, and why does this not contradict what you just proved?
\end{exercise}

\begin{exercise}[$\star\star$]\label{exo:b3:lambda-calculus:8}
Count the $\beta$-steps that normal order takes to reduce $\mathit{mult}\,c_m\,c_n$ to \omterm{def:b3:lambda-calculus:beta}{normal form}, as a function of $m$ and $n$, and explain why $n$ does not appear.
\end{exercise}

\begin{exercise}[$\star\star$]\label{exo:b3:lambda-calculus:9}
Using only the small-step rules of \code{if} and their congruence rule, show that $\mathtt{if}\ e_1\ \mathtt{then}\ e_2\ \mathtt{else}\ e_3 \to^* v$ holds exactly when $e_1 \to^* \mathtt{true}$ and $e_2 \to^* v$, or $e_1 \to^* \mathtt{false}$ and $e_3 \to^* v$. Deduce the big-step rules of \code{if}.
\end{exercise}

\begin{exercise}[$\star\star\star$]\label{exo:b3:lambda-calculus:10}
On the set $\{a, b, c, d\}$, let $b \to a$, $b \to c$, $c \to b$ and $c \to d$. Show that this system is locally \omterm{def:b3:lambda-calculus:rewriting}{confluent} but not \omterm{def:b3:lambda-calculus:rewriting}{confluent}, and say which hypothesis of Newman's lemma fails. Then prove that the measure used for the rules of $\mathit{add}$ in this chapter decreases when a rule is applied inside a larger term, not only at the top.
\end{exercise}

\begin{exercise}[$\star\star\star$]\label{exo:b3:lambda-calculus:11}
Prove that $\Omega$ has no \omterm{def:b3:lambda-calculus:beta}{normal form}. Show that $K\,I\,\Omega$, with $K = \lambda x.\,\lambda y.\,x$ and $I = \lambda z.\,z$, has one, that normal order reaches it in two steps, and that \omterm{def:b3:lambda-calculus:strategy}{call by value} never does.
\end{exercise}

\begin{exercise}[$\star\star\star$]\label{exo:b3:lambda-calculus:12}
The Krivine machine evaluates by \omterm{def:b3:lambda-calculus:strategy}{call by name} without substitution. Its state is a term, an environment (a list of closures, one per enclosing $\lambda$) and a stack of closures (pending arguments), where a closure is a term with its environment. An \omterm{def:b3:lambda-calculus:terms}{application} pushes its argument, as a closure, and continues with the function; an \omterm{def:b3:lambda-calculus:terms}{abstraction} pops a closure into its environment and continues with its body; a variable $k$ continues with the $k$-th closure of the environment. Implement it on \omterm{def:b3:lambda-calculus:debruijn}{de Bruijn} terms, with a function that reads a closure back as a term, and test that it reaches the same term as \omterm{def:b3:lambda-calculus:strategy}{call by name}, in the same number of $\beta$-steps.
\end{exercise}

\section{Weekend problem}\label{sec:b3:lambda-calculus:problem}

\begin{problem}[{Weekend problem --- Arithmetic with nothing but functions}]\label{pb:b3:lambda-calculus:1}
\omterm{def:b3:lambda-calculus:church}{Church numerals} are a complete arithmetic: every computable function on numbers can be written with them. This problem measures what it costs, counting $\beta$-steps to \omterm{def:b3:lambda-calculus:beta}{normal form}, under normal order (the strategy that always finds a \omterm{def:b3:lambda-calculus:beta}{normal form}) and applicative order. The program \code{pb_counts.rs} prints the counts it needs.

\omoutput{code/bachelor-3/05-lambda-calculus/outputs/pb_counts.txt}{$\beta$-steps to \omterm{def:b3:lambda-calculus:beta}{normal form} under normal order (\code{nor}) and applicative order (\code{app}); \code{exp c2 c0} is left out, its \omterm{def:b3:lambda-calculus:beta}{normal form} \texttt{\textbackslash{}x.\ x} not being a numeral.}\label{out:b3:lambda-calculus:pb}

\medskip\noindent\textbf{Part I --- Numerals and addition.}
\begin{enumerate}
\item Write $c_3$, and reduce $\mathit{succ}\,c_2$ to \omterm{def:b3:lambda-calculus:beta}{normal form} by hand, counting the steps.
\item Reduce $\mathit{plus}\,c_2\,c_3$ by normal order. Why do the counts of \cref{out:b3:lambda-calculus:church} for $\mathit{plus}$ not depend on the numbers?
\item Is the number of $\beta$-steps a fair measure of the work? What does each step of $\mathit{plus}\,c_m\,c_n$ copy, and how large is the result?
\item How many symbols does $c_n$ have? Compare with the binary representation of $n$.
\item Show that $\lambda m.\,\lambda n.\,m\ \mathit{succ}\ n$ also computes $m + n$.
\item Why does that second addition need more steps than $\mathit{plus}$ as $m$ grows?
\end{enumerate}

\medskip\noindent\textbf{Part II --- Multiplication and exponentiation.}
\begin{enumerate}[resume]
\item From \cref{out:b3:lambda-calculus:pb}, give the counts for $\mathit{mult}\,c_n\,c_3$ as a formula, and check it on two lines.
\item Give the counts for $\mathit{exp}\,c_2\,c_n$ under both strategies as formulas.
\item Explain the doubling under normal order: what does normal order copy that applicative order has already reduced?
\item Is normal order always slower? Give a term for which it takes fewer steps than applicative order.
\item What size does the \omterm{def:b3:lambda-calculus:beta}{normal form} of $\mathit{exp}\,c_2\,c_n$ have? Does its step count, under either strategy, measure the work of writing it out?
\item Which strategy would you choose for a calculator built on \omterm{def:b3:lambda-calculus:church}{Church numerals}, and what would call by need change?
\end{enumerate}

\medskip\noindent\textbf{Part III --- The predecessor.}
\begin{enumerate}[resume]
\item Show that after $k \ge 1$ \omterm{def:b3:lambda-calculus:terms}{applications} of the step $(a, b) \mapsto (b, b + 1)$ to $(0, 0)$, the pair is $(k - 1, k)$.
\item Why does $\mathit{pred}\,c_0$ give $c_0$?
\item From \cref{out:b3:lambda-calculus:pb}, give the normal-order count of $\mathit{pred}\,c_n$ for $n \ge 1$ as a formula.
\item How many steps does each round of the step take under normal order? Where do the remaining ones go?
\item Give the applicative-order count for $n \ge 2$ as a formula. Which strategy wins for large $n$?
\item Compare with the cost of a predecessor on binary numbers. What has the encoding cost?
\end{enumerate}

\medskip\noindent\textbf{Part IV --- Subtraction.}
\begin{enumerate}[resume]
\item Show that $\mathit{sub}\,c_m\,c_n = c_n\ \mathit{pred}\ c_m$ reduces to $c_{\max(m - n, 0)}$.
\item From the column \code{sub 10 n}, compute the differences between consecutive normal-order counts. What do you notice?
\item Suppose the $k$-th predecessor, for $k = 1, \dots, n$, costs as much as $\mathit{pred}$ applied to the numeral $c_{m - k + 1}$, plus 4 steps in all to unfold $\mathit{sub}$. Derive a formula for the count when $1 \le n \le m$, and check it against the table.
\item Under normal order, the predecessors are applied to terms that are not yet numerals. Why does the assumption of the previous question still hold?
\item What is the order of growth of the count in $m$ and $n$? At $m = n$?
\item Sum up: give the cost of $\mathit{pred}$ and of $\mathit{sub}$ under normal order, and say what a machine integer would cost instead.
\end{enumerate}
\end{problem}
